\subsection*{第2章の索引用語}

\begin{description}
  \item[\term{アセンブリ}] 機械語命令へ人が読める短い名前を付けた表記です。
  \item[\term{分岐}] 条件または無条件で次に実行する命令位置を変える処理です。
  \item[\term{ソースコード}] 人が読み書きするプログラミング言語で記した変換前のプログラムです。
  \item[\term{実行可能形式}] OSがロードして実行を開始できるファイル形式です。
  \item[\term{ローダ}] 実行可能形式をメモリへ配置し、実行準備を行う処理です。
  \item[\term{正規表現}] 文字列のパターンを記述する形式です。
  \item[\term{EBNF}] Extended Backus-Naur Formの略で、文法規則を記述する記法です。
  \item[\term{パーサ}] トークン列を文法に従って構造へ変換する処理です。
  \item[\term{ベースケース}] 再帰呼び出しを終了して値を返す条件です。
  \item[\term{前順走査}] 親ノードを子ノードより先に訪れる木の走査です。
  \item[\term{中順走査}] 二分木で左の子、親、右の子の順に訪れる走査です。
  \item[\term{後順走査}] 子ノードを親ノードより先に訪れる木の走査です。
  \item[\term{優先順位}] 演算子がどの順に結合するかを決める規則です。
  \item[\term{結合規則}] 同じ優先順位の演算子を左または右のどちらからまとめるか決める規則です。
  \item[\term{意味解析}] 構文が表す名前、型、作用域などの意味を確認する処理です。
\end{description}
